%%%%%%%%%%%%%%%%%%%%%%%%%%%%%%%%%%%%%%%%%%%%%%%%%%%%%%%%%%%%%%%%%%%%%
%% PolyJarvis: LLM Agent for Autonomous Polymer MD Simulations
%% Authors: Alexander Zhao, Achuth Chandrasekhar, Amir Barati Farimani

%%%%%%%%%%%%%%%%%%%%%%%%%%%%%%%%%%%%%%%%%%%%%%%%%%%%%%%%%%%%%%%%%%%%%
\documentclass[journal=acsami,manuscript=article]{achemso}

%%%%%%%%%%%%%%%%%%%%%%%%%%%%%%%%%%%%%%%%%%%%%%%%%%%%%%%%%%%%%%%%%%%%%
%% Packages
%%%%%%%%%%%%%%%%%%%%%%%%%%%%%%%%%%%%%%%%%%%%%%%%%%%%%%%%%%%%%%%%%%%%%
\usepackage[version=3]{mhchem}
\usepackage{multirow}
\usepackage{makecell}
\usepackage{booktabs}
\usepackage{amsmath}
\usepackage{amssymb}
\usepackage{graphicx}
\usepackage{xcolor}
\usepackage{algorithm}
\usepackage{algpseudocode}
\usepackage{hyperref}
\usepackage{float}

%% Placeholder command for items the user must fill in
\newcommand{\placeholder}[1]{\textcolor{red}{\textbf{[PLACEHOLDER: #1]}}}
%% Revision-change marker. MARKED version: highlights every revision change in blue.
%% The clean version (main_clean.tex) defines this as a no-op. ACS "marked manuscript".
\newcommand{\marked}[1]{\textcolor{blue}{#1}}
%% Scaffold block: a labeled, boxed bullet-note marking a planned revision change.
%% Bullet scaffolding only this pass — replaced by finished prose in the follow-up pass.
%% #1 = reviewer-comment tag (e.g. R1M4); #2 = itemize body.
\newcommand{\scaffold}[2]{\par\smallskip\noindent\marked{\fbox{\parbox{0.95\linewidth}{%
\footnotesize\textbf{[SCAFFOLD --- #1]}\par #2}}}\par\smallskip}

%%%%%%%%%%%%%%%%%%%%%%%%%%%%%%%%%%%%%%%%%%%%%%%%%%%%%%%%%%%%%%%%%%%%%
%% Meta-data
%%%%%%%%%%%%%%%%%%%%%%%%%%%%%%%%%%%%%%%%%%%%%%%%%%%%%%%%%%%%%%%%%%%%%
\author{Alexander Zhao}
% \email{arz2@andrew.cmu.edu}
\affiliation{Department of Materials Science and Engineering, Carnegie Mellon University, 5000 Forbes Avenue, Pittsburgh, PA 15213, USA}

\author{Achuth Chandrasekhar}
\affiliation{Department of Mechanical Engineering, Carnegie Mellon University, 5000 Forbes Avenue, Pittsburgh, PA 15213, USA}

\author{Amir Barati Farimani}
\email{barati@cmu.edu}
\affiliation{Department of Mechanical Engineering, Carnegie Mellon University, 5000 Forbes Avenue, Pittsburgh, PA 15213, USA}

%%%%%%%%%%%%%%%%%%%%%%%%%%%%%%%%%%%%%%%%%%%%%%%%%%%%%%%%%%%%%%%%%%%%%
%% Title
%%%%%%%%%%%%%%%%%%%%%%%%%%%%%%%%%%%%%%%%%%%%%%%%%%%%%%%%%%%%%%%%%%%%%
% [R1 Minor 1] Title — scope-accurate revision (SCAFFOLD; preamble, so kept as a comment note).
%   Reviewer: "autonomous" / "all-atom MD of polymers" overclaims the demonstrated scope.
%   - Make the title scope-accurate to amorphous homopolymers and an orchestrated agent.
%   - Candidate: "PolyJarvis: An LLM-Orchestrated Agent for Automated All-Atom Molecular
%     Dynamics of Amorphous Homopolymers."
%   - Decision deferred to final-prose pass once the 9-system benchmark scope is locked.
%   (Also surfaced as a bullet in the visible abstract scaffold box below.)
\title{PolyJarvis: Autonomous Large Language Model Agent for All-Atom Molecular Dynamics Simulation of Polymers}

\keywords{large language models, molecular dynamics, polymer simulation, autonomous agents, glass transition temperature, Model Context Protocol}

%%%%%%%%%%%%%%%%%%%%%%%%%%%%%%%%%%%%%%%%%%%%%%%%%%%%%%%%%%%%%%%%%%%%%
\begin{document}

%%%%%%%%%%%%%%%%%%%%%%%%%%%%%%%%%%%%%%%%%%%%%%%%%%%%%%%%%%%%%%%%%%%%%
%% Abstract
%%%%%%%%%%%%%%%%%%%%%%%%%%%%%%%%%%%%%%%%%%%%%%%%%%%%%%%%%%%%%%%%%%%%%
\begin{abstract}
All-atom molecular dynamics (MD) simulations can predict polymer properties from molecular structure, yet their execution requires specialized expertise in force field selection, system construction, equilibration, and property extraction.
We present PolyJarvis, an agent that couples a large language model (LLM) with the RadonPy simulation platform through Model Context Protocol (MCP) servers, enabling end-to-end polymer property prediction from natural language input.
Given a polymer name or SMILES string, PolyJarvis autonomously executes monomer construction, charge assignment, polymerization, force field parameterization, GPU-accelerated equilibration, and property calculation.
Validation is conducted on polyethylene (PE), atactic polystyrene (aPS), poly(methyl methacrylate) (PMMA), and poly(ethylene glycol) (PEG). Results show density predictions within 0.1--4.8\% and bulk moduli within 17--24\% of reference values for aPS and PMMA.
PMMA glass transition temperature ($T_g$) (395~K) matches experiment within +10--18~K, while aPS and PEG overestimate $T_g$ by +43--47~K (vs.\ upper experimental bounds).
Of the 8 property--polymer combinations with directly comparable experimental references, 5 meet strict acceptance criteria. For cases lacking suitable amorphous-phase experimental, agreement with prior MD literature is reported separately. The remaining $T_g$ failures are attributable primarily to the intrinsic MD cooling-rate bias rather than agent error.
This work demonstrates that LLM-driven agents can autonomously execute polymer MD workflows producing results consistent with expert-run simulations.
\scaffold{R1 Minor 1+2 --- Title \& Abstract}{%
\textbf{Title (R1 Minor 1):} make scope-accurate, e.g.\ ``PolyJarvis: An LLM-Orchestrated Agent for Automated All-Atom Molecular Dynamics of Amorphous Homopolymers.''\par
\textbf{Abstract (R1 Minor 2):} restructure (final prose pass; numbers \placeholder{pending 9-system $\times$ 5-replicate runs}):
\begin{itemize}
\item Lead with the \emph{strict} (experiment-only) scorecard: state $X/Y$ pass rate up front.
\item State plainly that PE density is overestimated ($\sim$25\%) and that two of the strict-benchmark $T_g$ cases fail the $\pm$20~K criterion.
\item Report replicate \emph{mean $\pm$ std} as the headline statistic (not best replicate).
\item Note PEG density/$K$ are validated against MD literature, not experiment.
\item Remove phrasing implying broad quantitative accuracy; soften ``autonomously execute.''
\item Tag as proof-of-concept across 9 amorphous homopolymers spanning 7 chemistries.
\end{itemize}}
\end{abstract}

%%%%%%%%%%%%%%%%%%%%%%%%%%%%%%%%%%%%%%%%%%%%%%%%%%%%%%%%%%%%%%%%%%%%%
%% Introduction
%%%%%%%%%%%%%%%%%%%%%%%%%%%%%%%%%%%%%%%%%%%%%%%%%%%%%%%%%%%%%%%%%%%%%
\section{Introduction}

Computational prediction of polymer properties such as $T_g$, density, and mechanical moduli is essential for accelerating materials discovery.\cite{bicerano2002prediction}
All-atom MD simulations provide a first-principles route to these predictions,\cite{kremer2009multiscale} but their practical adoption is hindered by the deep expertise required across force field selection, system construction, multi-stage equilibration, and property-specific analysis.\cite{larsen2011molecular, hayashi2022radonpy}
The manual, decision-intensive nature of these workflows limits reproducibility, with studies on the same system frequently reporting divergent results.\cite{patrone2016uncertainty, suter2025ensemble}

Significant progress has been made in automating the polymer simulation pipeline.
RadonPy\cite{hayashi2022radonpy} provides an end-to-end framework calculating 15 properties for more than 1000 polymers. SPACIER\cite{nanjo2025spacier} extends the framework with Bayesian optimization. Other notable tools like Polymatic,\cite{abbott2013polymatic} and Polyply\cite{grunewald2022polyply} automate specific MD stages. The PolyArena benchmark\cite{polyarena2025} establishes a standardized evaluation of MD-predicted properties.
Machine learning approaches offer rapid prediction but sacrifice physical interpretability and require extensive training data\cite{tao2021benchmarking, chaudhari2024alloybert}. However, existing tools execute predefined protocols without making informed decisions based on the specific polymer system.

Agents based on large language models (LLMs) have opened new possibilities for scientific automation.
For example, ChemCrow\cite{bran2024chemcrow} integrated chemistry tools with GPT-4 for organic synthesis, MDCrow\cite{campbell2025mdcrow} completed 72\% of 25 biomolecular MD tasks, DynaMate\cite{guilbert2025dynamate} introduced self-correcting protein--ligand MD, and an autonomous simulation agent was explored for polymer conformations.\cite{liu2025asa}
In materials science specifically, modular LLM agents have been applied to multi-task computational workflows,\cite{chaudhari2025modular} catalyst discovery,\cite{ock2026adsorb} materials knowledge navigation,\cite{zhuang2026natural} polymer design,\cite{nigam2026polymeragent} and biomolecular MD automation.\cite{chandrasekhar2025namdagent}
However, few have demonstrated an LLM agent autonomously executing validated all-atom polymer MD simulations with quantitative experimental benchmarking (Table~\ref{tab:comparison_agents}).

\begin{table}[ht]
\centering
\caption{Comparison of LLM-based MD automation frameworks.}
\label{tab:comparison_agents}
\small
\begin{tabular}{@{}llll@{}}
\toprule
\textbf{System} & \textbf{Domain} & \textbf{Validation} & \textbf{Autonomy} \\
\midrule
ChemCrow\cite{bran2024chemcrow} & Organic synth. & Expert eval. & Full \\
MDCrow\cite{campbell2025mdcrow} & Biomolecular & Task completion & Full \\
DynaMate\cite{guilbert2025dynamate} & Protein--ligand & Task completion & Full \\
MDAgent\cite{shi2025mdagent} & Materials MD & Expert eval. & Semi-auto \\
ASA\cite{liu2025asa} & Polymer conf. & Chain stats & Partial \\
Polymer-Agent\cite{nigam2026polymeragent} & Polymer design & Expert eval. & Full \\
NAMD-Agent\cite{chandrasekhar2025namdagent} & Protein MD & Thermo.\ params & Full \\
\textbf{PolyJarvis} & \textbf{Polymer MD} & \textbf{Exp.\ bench.} & \textbf{Full} \\
\bottomrule
\end{tabular}
\end{table}

\scaffold{R2 Comment 1 --- Framework comparison (Table 1)}{%
Add a paragraph after Table~\ref{tab:comparison_agents} comparing PolyJarvis to its closest prior frameworks:
\begin{itemize}
\item Contrast scope (polymer MD vs.\ small-molecule/protein domains of ChemCrow, MDCrow, DynaMate).
\item Contrast pipeline integration: PolyJarvis spans construction$\to$equilibration$\to$property extraction end-to-end vs.\ partial-stage tools.
\item Contrast validation scale/type: quantitative experimental benchmarking vs.\ task-completion or expert-eval metrics.
\item Be explicit where prior tools \emph{match or exceed} PolyJarvis (e.g.\ MDCrow task breadth; Polymer-Agent design focus) --- honest positioning.
\item Property coverage: $T_g$, density, $K$ across amorphous homopolymers.
\end{itemize}}
Here we present PolyJarvis, an LLM-powered autonomous agent for polymer property prediction through all-atom MD.
PolyJarvis integrates Anthropic's Claude\cite{anthropic2024claude} with RadonPy\cite{hayashi2022radonpy} via MCP\cite{anthropic2024mcp} servers.
The key distinction from prior work is \marked{LLM-orchestrated protocol adaptation across the workflow: classifying polymers to apply encoded force-field and charge defaults, applying convergence criteria, and composing recovery paths when stages fail. Most of these choices are deterministic lookups from structured tool interfaces; the agent's distinct contribution is integrating these encoded criteria with run-specific numerical context---most visibly in error recovery and for chemistries outside the validated table---rather than relying on hard-coded heuristics.}
We validate on four polymers spanning major backbone classes (PE, aPS, PMMA, PEG), predicting $T_g$, density, and bulk modulus with quantitative comparison to literature.


%%%%%%%%%%%%%%%%%%%%%%%%%%%%%%%%%%%%%%%%%%%%%%%%%%%%%%%%%%%%%%%%%%%%%
%% Methods
%%%%%%%%%%%%%%%%%%%%%%%%%%%%%%%%%%%%%%%%%%%%%%%%%%%%%%%%%%%%%%%%%%%%%
\section{Methods}

\subsection{System Architecture}

PolyJarvis operates through a client--server architecture mediated by MCP,\cite{anthropic2024mcp} comprising three components (Figure~\ref{fig:architecture}):
(1) an LLM backbone (Claude Sonnet 4.5)\cite{anthropic2024claude} that processes prompts, maintains workflow state, and makes simulation decisions;
(2) a local MCP server wrapping RadonPy\cite{hayashi2022radonpy} for molecular construction, force field assignment, and property analysis; and
(3) a remote MCP server managing GPU-accelerated LAMMPS\cite{thompson2022lammps} simulations on Lambda Labs cloud infrastructure (4$\times$ Quadro RTX 6000 GPUs).
The MCP framework provides a standardized interface through which the LLM discovers and invokes computational tools covering construction, simulation, analysis, and job monitoring (full tool inventory in Section~S1 of the Supporting Information (SI)).
\scaffold{R1 Minor 4 --- What MCP adds beyond API wrappers}{%
Add 1--2 sentences here (and expand in SI~S1) clarifying MCP's technical contribution:
\begin{itemize}
\item Typed tool schemas the LLM reads at runtime to reason about valid inputs (not hand-coded call sites).
\item Multi-turn tool invocation with persisted context across a long workflow; async job lifecycle (submit/poll/collect).
\item Standardized structured error propagation back into the reasoning loop (enables recovery).
\item Distinction from a workflow manager: \emph{no fixed DAG} --- the LLM determines execution order at runtime.
\end{itemize}}


\begin{figure}
\centering
\includegraphics[width=\textwidth]{figures/figure1_architecture.pdf}
\caption{PolyJarvis system architecture.
The LLM agent receives natural language polymer specifications and orchestrates the simulation workflow through two MCP servers.
The local server handles molecular construction using RadonPy and the remote server executes GPU-accelerated LAMMPS simulations.
Arrows indicate bidirectional communication via MCP.}
\label{fig:architecture}
\end{figure}


\subsection{Agent Decision-Making and Simulation Protocol}

A distinguishing feature of PolyJarvis is polymer-specific decision-making at each workflow stage.

\textbf{Polymer classification and force field selection.}
The agent classifies polymers against 21 backbone classes from the PoLyInfo scheme\cite{otsuka2011polyinfo} and selects force field, charge method, and electrostatics handling accordingly.

\textbf{System construction.}
The agent targets 10 polymer chains per system, consistent with RadonPy-validated minimum sampling.\cite{hayashi2022radonpy}
Chain lengths are adapted per system ($n = 100$--150 for PE, 62 for aPS, 100 for PMMA and PEG), yielding 6,020--15,020 atoms.
Amorphous cells are typically generated at 0.05~g/cm$^3$ initial density and compressed during equilibration.

\textbf{Equilibration.}
Staged protocols follow the compression/decompression method of Larsen et al.:\cite{larsen2011molecular} energy minimization, NVT heating, NPT compression at elevated pressure, decompression to 1~atm with full electrostatics, NPT cooling, and NVT production.
The agent adapts stage count, annealing cycles, and compression method based on system behavior.

\textbf{Property calculation.}
$T_g$ is extracted from stepwise-cooling density--temperature data via bilinear fitting, with sweep ranges and equilibration times adapted to each polymer's expected $T_g$.
Density is time-averaged over the final 50\% of NPT production at 300~K.
Bulk modulus is computed from NPT volume fluctuations: $K = k_BT\langle V\rangle / \langle(\delta V)^2\rangle$.\cite{allen2017computer}
Structural diagnostics (RDFs, end-to-end distances) validate physical reasonableness of simulated systems.
\\
A representative interaction of the agent is illustrated in Figure~\ref{fig:conversation}. A detailed decision-making protocol is provided in Section S2 of the SI, and property-extraction details are given in Sections S3--S4.

\begin{figure}
\centering
\includegraphics[width=0.85\columnwidth]{figures/figure2_conversation.pdf}
\caption{Representative PolyJarvis agent conversation.
The user submits a natural language task prompt specifying the polymer system, target properties, and computational resources.
The agent responds with a proposed workflow and any clarifying questions before proceeding autonomously.
After the simulation completes, the agent reports predicted property values alongside experimental references and percent error.}
\label{fig:conversation}
\end{figure}
\scaffold{R1 Minor 5 --- Real conversation trace (figure deferred)}{%
Replace the placeholder conversation figure with a real shortened trace:
\begin{itemize}
\item Use the PMMA4 D-05 convergence-gate-override incident as the primary trace.
\item Show: (1) \texttt{check\_equilibration\_comprehensive} call + structured return (tool output), (2) agent's per-gate physical argument (LLM reasoning), (3) the EXTEND-overridden Tier-3 decision.
\item Annotate green = tool output, blue = LLM reasoning, orange = ``Tier 3 --- runtime integration.''
\item Figure regeneration deferred; this pass scaffolds the caption/intent only.
\end{itemize}}
\clearpage
% [R1 Major 1] New "Decision Architecture" subsection — three-tier taxonomy + honest caveat.
\subsection{Decision Architecture}
\scaffold{R1 Major 1 --- Decision architecture (Tier 1/2/3)}{%
New subsection separating what is encoded from what the LLM infers at runtime. Add a condensed table (one row per tier):
\begin{itemize}
\item \textbf{Tier 1 --- encoded in tool schemas / system prompts} (LLM reads, does not reason): FF routing table, SMILES connectivity convention, convergence thresholds and gate-override conditionals, partial-sweep acceptance criteria, stage execution order.
\item \textbf{Tier 2 --- encoded structured parameters} (LLM fills template, no judgment): $T_\text{equil}$, DP, nchain, annealing cycles, density targets, pre-computed $T_\text{workflow}$.
\item \textbf{Tier 3 --- runtime integration of literature principles with run-specific numbers} (the LLM contribution): recovery-path composition for unanticipated failures, and best-judgment chemistry decisions for classes outside the validated table --- decisions not reducible to a pre-written conditional. Screen every candidate by one test: \emph{could the if-then have been written in advance?} If yes, it belongs in Tier 1, not here.
\item \textbf{Honest caveat (must state):} the physical principles are literature-derived and encoded in system context; the LLM does \emph{not} generate novel polymer physics. Its contribution is applying those principles to specific numerical runtime context --- a synthesis not reducible to a finite rule set.
\item \textbf{Generalizability (R1 Major 2 / Major 1):} the framework runs \emph{novel} systems not anchored by literature benchmarks, with the agent making best-judgment chemistry decisions (FF routing, SMILES/atom-typing fixes) per system --- evidence of genuine inference vs.\ a lookup table.
\item Representative Tier-3 incidents to cite (from run logs, populated in final pass): melt-input diagnosis + recovery, SMILES atom-typing correction, off-table chemistry decision --- each accompanied by the conversation trace showing it was not pre-scripted.
\end{itemize}}
\scaffold{R1 Major 1 / R1 Minor 3 --- Language audit}{%
Replace promotional ``intelligent decision-making'' phrasing with tier-appropriate language throughout (final pass):
\begin{itemize}
\item Tier 1/2 actions $\to$ ``the orchestrator selects from encoded parameters'' / ``the workflow executes the fixed protocol.''
\item Tier 3 actions $\to$ ``the agent integrates literature criteria with run-specific numerical context to determine\ldots''
\item Headline replacement: ``intelligent decision-making'' $\to$ ``LLM-integrated protocol adaptation.''
\item Audit \texttt{intelligent / smart / decides / autonomously} at: abstract; intro (``intelligent decision-making at every workflow stage''); Methods (``polymer-specific decision-making''); Results \S\ref{sec:agent_decisions}.
\end{itemize}}

\subsection{Benchmark Systems}

Four polymers were selected to span major backbone classes with extensive experimental characterization (Table~\ref{tab:benchmarks}).
Each system was run three times. Initial runs informed protocol refinements in subsequent runs (Table~\ref{tab:protocol_evolution}).

\begin{table*}[ht]
\centering
\caption{Benchmark polymer systems and reference properties.
Experimental (Exp.) values from the Polymer Handbook, 4th ed.\ (PH)\cite{brandrup1999polymerhandbook} and the PoLyInfo database (PI);\cite{otsuka2011polyinfo}
literature MD (MD Lit.) values where noted. A dash (---) indicates no suitable reference of that type for strict validation. Such cases are discussed separately against prior MD literature where appropriate.}
% [R1 Minor 6] Visually separate experimental vs MD-literature vs excluded references — see scaffold after table.
\label{tab:benchmarks}
\footnotesize
\begin{tabular}{ll cc cc cc}
\toprule
& & \multicolumn{2}{c}{\textbf{$T_g$ (K)}} & \multicolumn{2}{c}{\textbf{$\rho$ (g/cm$^3$)}} & \multicolumn{2}{c}{\textbf{$K$ (GPa)}} \\
\cmidrule(lr){3-4} \cmidrule(lr){5-6} \cmidrule(lr){7-8}
\textbf{Polymer} & \textbf{SMILES} & \textbf{Exp.} & \textbf{MD Lit.} & \textbf{Exp.} & \textbf{MD Lit.} & \textbf{Exp.} & \textbf{MD Lit.} \\
\midrule
PE   & \texttt{*CC*}              & 145--243\textsuperscript{a} & 200--327\textsuperscript{b} & 0.855\textsuperscript{c} & ---          & $\sim$1\textsuperscript{d} & 1--4\textsuperscript{e} \\
aPS  & \texttt{*CC(c1ccccc1)*}    & 353--373\textsuperscript{f} & 449--484\textsuperscript{g} & 1.04--1.065\textsuperscript{f} & --- & 3.55--4.5\textsuperscript{h} & --- \\
PMMA & \texttt{*CC(C)(C(=O)OC)*}  & 377--385\textsuperscript{f} & 415--513\textsuperscript{i} & 1.18--1.20\textsuperscript{f} & ---  & 4.1--4.8\textsuperscript{j} & --- \\
PEG  & \texttt{*CCO*}             & 206--213\textsuperscript{k} & 225--302\textsuperscript{l} & ---  & 1.10--1.13\textsuperscript{m} & --- & 1.2--3.5\textsuperscript{n} \\
\bottomrule
\multicolumn{8}{l}{\scriptsize \textsuperscript{a} PH: three conflicting values ($-128\pm5$, $-80\pm10$, $-30\pm15$~$^\circ$C); Boyer 1973.\cite{boyer1973pe}} \\
\multicolumn{8}{l}{\scriptsize \textsuperscript{b} Yang 2016\cite{yang2016pe} (200~K); Soldera 2006\cite{soldera2006tg} (285--327~K).} \\
\multicolumn{8}{l}{\scriptsize \textsuperscript{c} PH amorphous.\cite{brandrup1999polymerhandbook} \textsuperscript{d} PI: Aleman 1990. \textsuperscript{e} Morikami 1996;\cite{morikami1996bulkmodulus} 1--2 rubbery, 3--4 glassy.} \\
\multicolumn{8}{l}{\scriptsize \textsuperscript{f} PH;\cite{brandrup1999polymerhandbook} PI.\cite{otsuka2011polyinfo} \textsuperscript{g} Soldera 2006\cite{soldera2006tg}.} \\
\multicolumn{8}{l}{\scriptsize \textsuperscript{h} PH $\beta_T$;\cite{brandrup1999polymerhandbook} PI: 3.55~GPa. \textsuperscript{i} Tang 2022\cite{tang2022pmma} (481~K); Soldera 2006\cite{soldera2006tg} (415--513~K).} \\
\multicolumn{8}{l}{\scriptsize \textsuperscript{j} PH $\beta_T$;\cite{brandrup1999polymerhandbook} PI: 4.8~GPa. \textsuperscript{k} T\"{o}rm\"{a}l\"{a} 1974\cite{tormala1974pegtg} (213~K); PI: 206~K.} \\
\multicolumn{8}{l}{\scriptsize \textsuperscript{l} Wu 2011\cite{wu2011peo} (253--271~K); Wu 2019\cite{wu2019peo} (225~K).} \\
\multicolumn{8}{l}{\scriptsize \textsuperscript{m} Sponseller 2021;\cite{sponseller2021peg} Wu 2011.\cite{wu2011peo} \textsuperscript{n} Kacar 2018\cite{kacar2018peg} ($\sim$3.5); Wu 2019\cite{wu2019peo} ($\sim$1.2).} \\
\end{tabular}
\end{table*}
\scaffold{R1 Minor 6 --- Visually separate reference types}{%
Make the benchmark table's reference basis obvious at a glance (final pass + 5 new systems added):
\begin{itemize}
\item Use column/row shading or grouping to distinguish (A) strict experimental benchmark, (B) MD-literature context, (C) excluded / no comparable reference.
\item PE and PEG entries especially must show their benchmark basis without reading footnotes.
\item Add the 5 new polymers (PLA, PVC, PSU, cis-PBD, PEEK) as rows with their reference sources.
\end{itemize}}


%%%%%%%%%%%%%%%%%%%%%%%%%%%%%%%%%%%%%%%%%%%%%%%%%%%%%%%%%%%%%%%%%%%%%
%% Results and Discussion
%%%%%%%%%%%%%%%%%%%%%%%%%%%%%%%%%%%%%%%%%%%%%%%%%%%%%%%%%%%%%%%%%%%%%
\section{Results and Discussion}

Acceptance criteria for validation, consistent with established standards,\cite{afzal2021highthroughput, Harish2024, maicas2021epoxyMD} are: $|T_g^{\text{sim}} - T_g^{\text{exp}}| \leq 20$~K, density relative error $\leq$5\%, and bulk modulus relative error $\leq$30\%.
Strict validation is assessed only against experimental references. Literature MD values are used only as secondary contextual benchmarks when experimental amorphous-phase data are unavailable or not comparable to the present simulations.
\subsection{Agent Decision-Making}
\label{sec:agent_decisions}

The agent correctly classified all four polymers and selected force fields without human intervention.
Notable decisions include: employing GAFF2\cite{wang2004gaff} 
instead of GAFF2\_mod\cite{trag2019gaff2mod} for PEG given it is broadly validated, selecting Gasteiger charges for PE (negligible Coulombic contributions from CH$_2$ partial charges), removing PPPM during the aPS Run~2 $T_g$ sweep ($\sim$40\% speedup, reverted in Run~3), and switching PMMA from NPT barostat compression to \texttt{fix deform} after PPPM failure.

A key capability is iterative refinement across replicates (Table~\ref{tab:protocol_evolution}): the agent expanded aPS equilibration from 6 to 12 stages after Run~1, increased PE chain length from $n = 100$ to 150, and diagnosed a velocity re-initialization bug in PMMA.
Per-run parameters are provided in Section~S9 of the SI.

\begin{table}[ht]
\centering
\caption{Protocol refinements across replicate runs.}
\label{tab:protocol_evolution}
\small
\begin{tabular}{@{}llp{3.2cm}l@{}}
\toprule
\textbf{System} & \textbf{Run} & \textbf{Change} & \textbf{Impact} \\
\midrule
\multirow{2}{*}{PE} & 2 & $n$: 100$\to$150; 10 chains & Better chain stats \\
& 3 & 10~K step; 2~ns/step & $R^2 > 0.997$ \\
\midrule
\multirow{2}{*}{aPS} & 2 & 12-stage eq., 3$\times$ anneal & $T_g:498 \to 416K$  \\
& 3 & 8-stage, 1$\times$ anneal & $T_g = 466$~K \\
\midrule
PEG & 2 & 2~ns/step (from 1~ns) & $T_g = 275$~K \\
\midrule
\multirow{2}{*}{PMMA} & 2 & \texttt{fix deform} (vs.\ NPT) & Fixed PPPM crash \\
& 3 & Run~2 protocol, new seed & $T_g = 473$~K \\
\bottomrule
\end{tabular}
\end{table}

The agent autonomously resolved 13 simulation errors across four categories: PPPM out-of-range atoms (4), pair style mismatches (3), velocity re-initialization bugs (2), and $T_g$ extraction poor fits (4); full decision and error logs are provided in Section~S10 of the SI.
\scaffold{R1 Major 11 --- Error-recovery benchmark (fully pending)}{%
The ``13 errors'' count is provisional and \placeholder{will change once the 9-system revision runs complete}:
\begin{itemize}
\item Reframe as a reproducible benchmark: per-incident class, recovery action, \emph{pre-scripted vs.\ Tier-3 (LLM-inferred)}, \#retries (incl.\ failed attempts), outcome.
\item State how much human guidance (if any) was involved per recovery.
\item Full table assembled in the final pass from the completed revision-run logs --- not back-filled from prior exploratory runs.
\end{itemize}}


\subsection{Glass Transition Temperature}
\label{sec:tg_results}
\begin{figure*}[h]
\centering
\includegraphics[width=0.85\textwidth]{figures/figure3_Tg_curves.png}
\caption{Density versus temperature curves for the best replicate of each benchmark polymer, with bilinear fits (black lines), extracted $T_g$ (gold stars), and literature experimental ranges (green shaded area).}
\label{fig:tg_curves}
\end{figure*}
\scaffold{R1 Minor 8 --- Show $T_g$ replicates in main figure (deferred)}{%
Main $T_g$ figure currently shows only the best replicate. Replace (figure \placeholder{regenerated after 5-replicate runs}):
\begin{itemize}
\item Main panel: all replicates as scatter points with a mean\,$\pm$\,std bar per polymer (not best replicate).
\item Keep individual density--temperature curves for all replicates in the SI.
\item Caption to state the replicate count and that mean\,$\pm$\,std is the reported value.
\end{itemize}}

Table~\ref{tab:tg_comparison} presents all $T_g$ predictions; reference values are in Table~\ref{tab:benchmarks}.
Representative density--temperature curves are shown in Figure~\ref{fig:tg_curves}.
\scaffold{R1 Major 5 / R1 Major 6 --- Replicate statistics as primary (not best replicate)}{%
Current $T_g$ table reports per-run values with ``best replicate'' bolded; this risks optimistic selection. Revise (\placeholder{values pending 5-replicate runs}):
\begin{itemize}
\item Report \textbf{mean $\pm$ std over 5 replicates} as the \emph{primary} result; move best-replicate to secondary/SI discussion.
\item R1M5 sensitivity: demonstrate stability to fit method, cooling schedule, temperature grid, and equilibration duration (SI sensitivity table).
\item Show the full run-to-run spread (esp.\ aPS, PMMA) rather than only the favorable trajectory.
\item Expand the table to all 9 systems.
\end{itemize}}

\begin{table}[h]
\centering
\caption{Predicted $T_g$ (K) via F-statistic bilinear split.
Parenthetical values are $R^2$; bold = best replicate.
Experimental and MD literature ranges from Table~\ref{tab:benchmarks}.}
\label{tab:tg_comparison}
\small
\begin{tabular}{lcccc}
\toprule
\textbf{Polymer} & \textbf{Run 1} & \textbf{Run 2} & \textbf{Run 3} & \textbf{$\Delta T_g$\textsuperscript{a}} \\
\midrule
PE   & 291 (.998) & \textbf{281} (.999) & 304 (.997) & +38 \\
aPS  & 498 (.988) & \textbf{416} (.996) & 466 (.990) & +43 \\
PMMA & \textbf{395} (.998) & 489 (.996) & 473 (.988) & +10 \\
PEG  & 273 (.998) & 275 (1.00) & \textbf{260} (1.00) & +47 \\
\bottomrule
\multicolumn{5}{l}{\footnotesize \textsuperscript{a} Best replicate minus upper experimental bound.} \\
\end{tabular}
\end{table}

\textbf{PE.}
The best replicate (Run~2, $T_g = 281$K) is not included in the strict validation count because the experimental PE $T_g$ literature is not sufficiently well-defined to provide an unambiguous amorphous benchmark. Nevertheless, the predicted value falls within the broad 200--327K range reported in prior MD studies.\cite{yang2016pe,soldera2006tg}

\textbf{aPS.}
Run~2 ($T_g = 416$~K) overestimates the experimental range (353--373~K) by +43--63~K, in line with the +79~K systematic bias found by Afzal et al.\cite{afzal2021highthroughput} and the 449--484~K values of Soldera \& Metatla\cite{soldera2006tg} at comparable cooling rates.

\textbf{PMMA.}
Run~1 ($T_g = 395$~K, $R^2 = 0.998$) falls within +10--18~K of experiment (377--385~K), the only system satisfying the $\pm$20~K criterion and notably closer than published MD values of 415--513~K.\cite{tang2022pmma,soldera2006tg}

\textbf{PEG.}
Run~3 ($T_g = 260$~K) overestimates experiment (206--213~K\cite{tormala1974pegtg}) by +47~K, consistent with AA MD literature (253--271~K\cite{wu2011peo}).
All three replicates achieve $R^2 \geq 0.998$ with a narrow 15~K spread, the best inter-run consistency of the four systems.

Across the three polymers retained in the strict experimental $T_g$ benchmark (aPS, PMMA, and PEG), the MAE relative to experimental upper bounds is 33.3~K. This is comparable to the 20--33~K MAE reported by Suter et al.\cite{suter2025ensemble} with 10+ replica ensembles and better than the 81~K uncalibrated MAE of Afzal et al.\cite{afzal2021highthroughput}. PE is discussed separately because the experimental literature does not provide an unambiguous amorphous $T_g$ benchmark.
The dominant error source is the intrinsic cooling-rate limitation, not agent-introduced inaccuracies. Run-by-run fit statistics are tabulated in Section~S5 of the SI.
\scaffold{R1 Minor 11 --- Soften cooling-rate attribution}{%
Current text attributes $T_g$ failures solely to cooling-rate bias. Revise (final pass):
\begin{itemize}
\item ``consistent with cooling-rate bias, \emph{though force-field, chain-length, equilibration, and fitting contributions cannot be ruled out}.''
\item Drop ``not agent-introduced inaccuracies'' as an unqualified claim; the replicate spread (next pass) lets us speak quantitatively.
\item Mirror the same softening in the Abstract and Limitations.
\end{itemize}}
\scaffold{R1 Minor 14 --- MD agreement $\neq$ experimental validation}{%
Audit every ``consistent with prior MD'' sentence (here and aPS/PMMA/PEG paragraphs above) and add the qualifier:
\begin{itemize}
\item ``Agreement with prior MD indicates methodological consistency but does not independently validate against experiment; both may share the same systematic (e.g.\ cooling-rate) biases.''
\item Specifically tag the aPS/PEG comparisons where MD literature is the only reference.
\end{itemize}}


\subsection{Density}
\label{sec:density}

Equilibrium densities at 300~K are compared in Table~\ref{tab:density_comparison} and Figure~\ref{fig:density_parity}.

\begin{table}[H]
\centering
\caption{Predicted vs.\ reference densities at 300~K.
Simulated values are ensemble means $\pm$ 1~s.d.\ across three replicates from NPT production runs.}
\label{tab:density_comparison}
\begin{tabular}{lcccc}
\toprule
\textbf{Polymer} & \textbf{$\rho^{\text{sim}}$ (g/cm$^3$)} & \textbf{$\rho^{\text{ref}}$ (g/cm$^3$)} & \textbf{Error (\%)} & \textbf{Pass?} \\
\midrule
PE   & 1.069 $\pm$ 0.007 & 0.855\textsuperscript{a} & +25.0 & No \\
aPS  & 1.054 $\pm$ 0.002 & 1.053\textsuperscript{a} & +0.1 & \textbf{Yes} \\
PMMA & 1.124 $\pm$ 0.011 & 1.18\textsuperscript{a} & $-$4.8 & \textbf{Yes} \\
PEG  & 1.099 $\pm$ 0.004 & 1.10--1.13\textsuperscript{b} & $-$1 to 0 & \textbf{Yes} \\
\bottomrule
\multicolumn{5}{l}{\footnotesize \textsuperscript{a} Experimental (PH).\cite{brandrup1999polymerhandbook} \textsuperscript{b} MD amorphous-phase refs;\cite{sponseller2021peg,wu2011peo,kacar2018peg} no exp.\ amorphous PEO $\rho$ available.} \\
\end{tabular}
\end{table}

\begin{figure}[H]
\centering
\includegraphics[width=0.5\textwidth]{figures/figure4_density_parity.pdf}
\caption{Predicted versus experimental density at 300~K.
Large symbols show ensemble means $\pm$ 1 s.d.; small symbols show individual runs.
The shaded band marks $\pm$5\%.}
\label{fig:density_parity}
\end{figure}
\scaffold{R1 Minor 7 --- Make PE density failure visually prominent (deferred)}{%
\placeholder{Figure regenerated after PE re-run}:
\begin{itemize}
\item Render the PE point distinctly (red / hatching / annotation arrow) so the largest density error is not buried among the aPS/PMMA/PEG successes.
\item Call PE out directly in the caption --- do not soften.
\item If the corrected PE re-run improves density, show original vs.\ corrected.
\end{itemize}}

Three of four systems pass the $\pm$5\% criterion: aPS (1.054~g/cm$^3$, +0.1\%) falls within the PH range, PMMA ($-$4.8\%) is at the acceptance boundary, and PEG (1.099~g/cm$^3$) matches all-atom (AA) MD amorphous-phase references\cite{sponseller2021peg,wu2011peo}.
PE's +25\% overestimation is systematic across all runs, which is due to an agent-introduced protocol refinement error. The PE systems were driven into a kinetically trapped state, resulting in an unphysically high density.
\scaffold{R1 Major 4 / R2 Comment 2 --- PE density post-mortem + autonomous detection}{%
Expand into a proper post-mortem (corrected numbers \placeholder{pending PE re-run}):
\begin{itemize}
\item Root-cause analysis: identify the specific refinement that over-densified (compression schedule / barostat lag / cooling rate / box size); show evidence of kinetic trapping vs.\ relaxation (MSD plateau, $R_{ee}$).
\item Report corrected density from the extended-protocol re-run \emph{alongside} the original +25\% to show the fix.
\item Autonomous-detection discussion (R2): propose a heuristic the agent can apply --- e.g.\ flag when final $\rho$\,/\,estimated crystalline $\rho > 0.95$, or density plateau persistently above class target --- and state whether the current workflow has it or it is added.
\item Do NOT soften: PE is the largest error and a baseline polymer; keep it prominent.
\end{itemize}}


\subsection{Bulk Modulus}

Isothermal bulk moduli from NPT production runs at 300~K are presented in Table~\ref{tab:bulkmod_comparison} and Figure~\ref{fig:bulk_modulus}; volume statistics underlying these calculations are given in Section~S7 of the SI.

\begin{table}[H]
\centering
\caption{Predicted vs.\ reference isothermal bulk moduli at 300~K.
Simulated values are ensemble means $\pm$ 1~s.d.\ across three replicates.}
\label{tab:bulkmod_comparison}
\begin{tabular}{lcccc}
\toprule
\textbf{Polymer} & \textbf{$K^{\text{sim}}$ (GPa)} & \textbf{$K^{\text{ref}}$ (GPa)} & \textbf{Error (\%)} & \textbf{Pass?} \\
\midrule
PE   & 2.87 $\pm$ 0.31 & 1--4\textsuperscript{a}          & Within range & \textbf{Yes\textsuperscript{a}} \\
aPS  & 2.71 $\pm$ 0.48 & 3.55--4.5\textsuperscript{b}     & $-$24\textsuperscript{d} & \textbf{Yes} \\
PMMA & 3.42 $\pm$ 0.38 & 4.1--4.8\textsuperscript{b}      & $-$17\textsuperscript{d} & \textbf{Yes} \\
PEG  & 2.82 $\pm$ 0.20 & 1.2--3.5\textsuperscript{c}      & $-$19\textsuperscript{e} & \textbf{Yes} \\
\bottomrule
\multicolumn{5}{l}{\footnotesize \textsuperscript{a} MD lit.: Morikami 1996;\cite{morikami1996bulkmodulus} 1--2~GPa (rubbery), 3--4~GPa (glassy).} \\
\multicolumn{5}{l}{\footnotesize \textsuperscript{b} Experimental (PH/PI). }\\
\multicolumn{5}{l}{\footnotesize \textsuperscript{c} Coarse-grained (CG) MD: Wu 2019\cite{wu2019peo} (1.2~GPa); AA MD: Kacar 2018\cite{kacar2018peg} (3.5~GPa).} \\
\multicolumn{5}{l}{\footnotesize \textsuperscript{d} Vs.\ lower bound of experimental range. \textsuperscript{e} Vs.\ AA MD ref.\ (3.5~GPa); no exp.\ amorphous $K$ for PEO.} \\
\end{tabular}
\end{table}

\begin{figure}[H]
\centering
\includegraphics[width=0.55\textwidth]{figures/figure5_bulk_modulus.pdf}
\caption{Predicted versus reference bulk modulus at 300~K.}
\label{fig:bulk_modulus}
\end{figure}

\scaffold{R1 Major 9 --- Bulk-modulus statistical robustness}{%
NPT-fluctuation $K$ is sensitive to barostat, sampling length, and volume autocorrelation. Add to main text (details in SI~S7; \placeholder{barostat-sensitivity values pending}):
\begin{itemize}
\item Report volume autocorrelation time, effective sample size $N_\text{eff}$, and block-averaged uncertainty per run.
\item Report sensitivity to barostat relaxation time $\tau_P$ (varied $2\times$ in both directions); compare $K_\text{fluc}$ vs.\ $\tau_P$-insensitive $K_\text{Born}$.
\item If $K_\text{fluc}$ is fragile, note practical reliability limits and prefer the Born/NVT estimate.
\end{itemize}}
All four systems pass the $\pm$30\% criterion against their respective references.
For aPS ($-$24\%) and PMMA ($-$17\%), simulated values fall below experimental bounds, with aPS showing notable inter-replicate variance ($\pm$0.48~GPa) due to sensitivity of volume fluctuations to equilibration quality.
PE (2.87~GPa) falls within the MD range of 1--4~GPa\cite{morikami1996bulkmodulus}; the elevated $K$ is consistent with PE's over-dense packing.
PEG ($2.82 \pm 0.20$~GPa) is $-$19\% vs.\ the AA reference of Kacar\cite{kacar2018peg} ($\sim$3.5~GPa).


\subsection{Structural Validation and Energy Convergence}
\label{sec:structural}

\scaffold{R1 Major 10 --- Stronger structural diagnostics}{%
RDF-to-unity and $R_{ee}$ are weak equilibration evidence (esp.\ below $T_g$ / for PE). Add (mostly auto-produced by the updated equil-checker; \placeholder{values pending revision runs}):
\begin{itemize}
\item $R_g$ distributions, chain MSD plateau check, density homogeneity (block density over final trajectory).
\item $C(t)$ end-to-end autocorrelation with KWW fit ($\tau_\text{relax}$) as the chain-relaxation indicator.
\item For PE specifically: explicit chain-relaxation vs.\ kinetic-trapping evidence (the point R1M4 hinges on).
\item Check for residual orientation / artificial packing.
\end{itemize}}
Figure~\ref{fig:rdf} shows representative C--C radial distribution functions for the best replicate of each polymer; per-system RDFs with all three replicate runs overlaid are provided in Section~S8 of the SI.
All RDFs exhibit sharp first-neighbor bond peaks at $\sim$1.53~\AA{}, well-resolved second- and third-neighbor peaks, and smooth convergence to $g(r) = 1$ within 0.5\% by $r = 15$~\AA{}, confirming homogeneous amorphous packing with no residual crystalline order.

\begin{figure}[H]
\centering
\includegraphics[width=0.95\columnwidth]{figures/figure6_rdf.pdf}
\caption{Carbon--carbon radial distribution functions for the best replicate of each benchmark polymer at 300~K (101-frame NVT trajectories, MDAnalysis InterRDF, $r_{\max} = 15$~\AA{}, 150 bins).
All four systems show clean amorphous structure with bond peaks decaying and $g(r) \to 1$ at long range.}
\label{fig:rdf}
\end{figure}

Table~\ref{tab:e2e} summarizes end-to-end distances across all runs.
For PE and aPS, 300~K falls below their simulated $T_g$, so chains are conformationally frozen and $\langle R \rangle$ values reflect quenched configurations rather than equilibrium statistics.
PEG shows the most consistent inter-run $\langle R \rangle$ (28--42~\AA{}, CV = 20\%), while PE exhibits large variation between Run~1 ($n = 100$, $\langle R \rangle = 35$~\AA{}) and Runs~2--3 ($n = 150$, $\langle R \rangle = 68$--75~\AA{}), directly reflecting the agent's chain-length refinement.

\begin{table}[H]
\centering
\caption{End-to-end distances at 300~K (10-chain ensemble means over 101 NVT frames).}
\label{tab:e2e}
\small
\begin{tabular}{@{}lcccc@{}}
\toprule
\textbf{Polymer} & \textbf{$n$} & \textbf{Run} & \textbf{$\langle R \rangle$ (\AA{})} & \textbf{$\langle R^2 \rangle$ (\AA{}$^2$)} \\
\midrule
\multirow{3}{*}{PE}   & 100  & 1 & 35.5 & 1439 \\
                       & 150 & \textbf{2} & 67.8 & 5374 \\
                       & 150 & 3 & 74.9 & 6262 \\
\midrule
\multirow{3}{*}{aPS}  & 62  & 1 & 24.7 & 741 \\
                       & 62  & \textbf{2} & 28.0 & 937 \\
                       & 62  & 3 & 42.0 & 2073 \\
\midrule
\multirow{3}{*}{PMMA} & 100 & \textbf{1} & 29.0 & 947 \\
                       & 100 & 2 & 41.7 & 1877 \\
                       & 100 & 3 & 42.8 & 1909 \\
\midrule
\multirow{3}{*}{PEG}  & 100 & 1 & 38.3 & 1777 \\
                       & 100 & 2 & 28.0 & 902 \\
                       & 100 & \textbf{3} & 41.6 & 1962 \\
\bottomrule
\end{tabular}
\end{table}

11/12 production runs, with the exception of PE run 1, achieved energy drift $<$0.5\% and density CV $<$2\%, confirming adequate equilibration across systems (complete convergence metrics in Section~S11 of the SI).
\scaffold{R1 Minor 13 --- Computational cost summary (main text)}{%
Add a short main-text summary (full table in SI~S12; \placeholder{values pending revision runs}):
\begin{itemize}
\item Wall time and GPU-hours per polymer; total GPU-hours for the campaign.
\item LLM API-call count and approximate cost.
\item Failure/retry overhead --- supports the practical-usability argument.
\end{itemize}}


\subsection{Validation Summary and Limitations}

\begin{table*}[ht]
\centering
\caption{Comprehensive validation summary.
$\checkmark$ = meets criterion against a matching experimental reference; $\times$ = fails against a  experimental reference; --- = excluded from the strict validation count because no directly comparable experimental reference is available.
$T_g$ values are the best replicate.
$T_g$ $\Delta$ computed vs.\ upper bound of experimental range.}
\label{tab:validation_summary}
\begin{tabular}{lcccc}
\toprule
\textbf{Property (Criterion)} & \textbf{PE} & \textbf{aPS} & \textbf{PMMA} & \textbf{PEG} \\
\midrule
$T_g$ ($\pm$20~K) & --- & $\times$ (+43~K) & $\checkmark$ (+10~K) & $\times$ (+47~K) \\
Density ($\pm$5\%) & $\times$ (+25.0\%) & $\checkmark$ (+0.1\%) & $\checkmark$ ($-$4.8\%) & --- \\
Bulk modulus ($\pm$30\%) & --- & $\checkmark$ ($-$24\%) & $\checkmark$ ($-$17\%) & --- \\
\bottomrule


\end{tabular}
\end{table*}
\scaffold{R1 Major 3 --- Strict vs.\ inclusive scorecard}{%
Replace the single ``5 of 8'' framing with two explicit scorecards (counts \placeholder{pending 9-system runs}):
\begin{itemize}
\item \textbf{Strict} row: pass/fail against \emph{experimental} references only; PE $T_g$ stated as \emph{excluded} (ambiguous experimental reference), not counted as a pass.
\item \textbf{Inclusive} row: also accepting MD-literature references (PEG density/$K$, etc.), clearly labelled as weaker evidence.
\item Present both as a two-row table so the reader sees the spread; do not let the inclusive number stand alone.
\item State PE density failure and the two failing strict $T_g$ cases explicitly in the surrounding text.
\end{itemize}}
Table~\ref{tab:validation_summary} summarizes the validation results. Using only comparable experimental benchmarks, 5 of 8 property--polymer combinations meet the predefined acceptance criteria. Cases lacking suitable experimental amorphous-phase references are excluded from the strict validation count and are instead compared separately against prior MD literature for contextual support.
Of the three polymers retained in the strict experimental $T_g$ benchmark, two fail the $\pm$20K criterion by +43 to +47K (vs.\ upper experimental bounds), consistent with the well-documented MD cooling-rate artifact.\cite{soldera2006tg}
PMMA ($T_g = 395$~K, $+$10--18~K) is the sole pass.

\scaffold{R1 Major 2 --- Expand / scope the validation set}{%
Four polymers (three textbook) is too small for the breadth of claims. Address by (\placeholder{results pending 9-system runs}):
\begin{itemize}
\item Expand to 9 systems spanning 7 chemistries (PE, aPS, PMMA, PEG + PLA, PVC, PSU, cis-PBD, PEEK): aliphatic, acrylate, ester, vinyl halide, sulfone, diene, aromatic ketone; 2 force fields; exp.\ $T_g$ range 181--463~K.
\item Report failures honestly --- failed pilot systems stay in the paper as failures, not replaced with easier ones.
\item Add an explicit \textbf{scope-exclusion} paragraph: amorphous cells only (no semicrystalline), no charged/ionomer systems, no copolymers.
\item Temper any remaining generality claims to the demonstrated scope.
\end{itemize}}
Key limitations:
(1) Validation is restricted to four amorphous homopolymers.
(2) PE density is overestimated by +25\% systematically, which is a failure of the agent's equilibration protocol.
(3) Run-to-run $T_g$ variability (aPS: 416--498~K; PMMA: 395--489~K) shows that equilibration quality can dominate prediction error.\cite{larsen2011molecular, suter2025ensemble}
(4) Limited experimental $K$ data for PE and PEG.
(5) MD cooling rates systematically overestimate $T_g$.\cite{soldera2006tg}
(6) LLM context window limits and session timeouts occasionally require user intervention.

% [R2 Comment 3] New subsection categorizing error origins.
\subsection{Sources of Prediction Error}
\scaffold{R2 Comment 3 --- Sources of prediction error}{%
New subsection categorizing each failed/borderline case by origin (anchor to FF-sensitivity results when they land):
\begin{itemize}
\item \textbf{Agent/protocol choice:} PE density (kinetically trapped over-densification).
\item \textbf{Force-field limitation:} \placeholder{candidate $T_g$/$K$ offsets pending FF sensitivity test}.
\item \textbf{Cooling-rate artifact:} systematic $T_g$ overestimation (aPS, PEG, PMMA replicates).
\item \textbf{Sampling insufficiency:} inter-replicate $K$ variance; equilibration-quality sensitivity.
\item Make clear which errors the autonomous workflow can vs.\ cannot currently self-detect.
\end{itemize}}
\scaffold{R1 Major 7 --- Baseline comparison (deferred)}{%
New subsection, \placeholder{pending stock-RadonPy baseline runs}:
\begin{itemize}
\item Run the benchmark polymers with stock RadonPy defaults (no LLM, no custom protocol); report density/$T_g$ side-by-side.
\item Frame PolyJarvis value as reduced user effort to reach the same-or-better result --- not necessarily higher accuracy.
\item Acknowledge honestly where RadonPy defaults already match.
\end{itemize}}
\scaffold{R1 Major 8 --- Force-field sensitivity (deferred)}{%
New subsection, \placeholder{pending FF sensitivity runs}:
\begin{itemize}
\item Pick one failed case (PE density or aPS $T_g$); run alternative FF/charge models (e.g.\ GAFF2+Gasteiger vs.\ GAFF2+AM1-BCC; OPLS-AA if feasible).
\item Report how much error shifts --- a negative result is still useful evidence.
\item Add a per-polymer FF-selection rationale paragraph.
\end{itemize}}


%%%%%%%%%%%%%%%%%%%%%%%%%%%%%%%%%%%%%%%%%%%%%%%%%%%%%%%%%%%%%%%%%%%%%
%% Conclusion
%%%%%%%%%%%%%%%%%%%%%%%%%%%%%%%%%%%%%%%%%%%%%%%%%%%%%%%%%%%%%%%%%%%%%
\section{Conclusions}

\scaffold{R1 Minor 15 --- Restrained conclusion tone}{%
Rewrite (final pass; \placeholder{scorecard numbers pending 9-system runs}):
\begin{itemize}
\item Lead with \emph{proof of concept}; state demonstrated scope (9 amorphous homopolymers, 7 chemistries) and identified failure modes (PE over-densification, cooling-rate $T_g$ bias).
\item Remove any phrasing implying a production-ready / generally-reliable autonomous platform.
\item Emphasize the need for broader validation (semicrystalline, copolymers, charged systems) as future work.
\item Keep the honest LLM-contribution framing (Tier-3 synthesis, not novel physics generation).
\end{itemize}}
We have presented PolyJarvis, an autonomous LLM agent that couples Claude with RadonPy and LAMMPS through MCP servers to execute complete polymer MD workflows from natural language input. Overall, the agent demonstrated iterative protocol refinement, error recovery, and root-cause analysis of force field limitations. Validation on four benchmark polymers shows that 5 of 8 property--polymer combinations with directly comparable experimental references meet strict quantitative acceptance criteria: aPS and PMMA bulk moduli and densities pass, and PMMA $T_g$ (395~K) matches experiment within +10--18~K. Of the three polymers retained in the strict experimental $T_g$ benchmark, two overestimate experiment by +43 to +47~K (best replicates vs.\ upper experimental bounds), consistent with the well-documented MD cooling-rate artifact\cite{soldera2006tg} and comparable to published all-atom MD results \cite{afzal2021highthroughput,suter2025ensemble}. A key system limitation remains: the agent failed to adapt its compression protocol to resolve kinetically trapped, over-dense states in PE.

This work demonstrates that an LLM agent can autonomously drive end-to-end polymer MD simulations, offering a shift in methodology where the agent handles workflow design and execution while manual effort is freed for ideation and experimental design. Future work will expand force field coverage (OPLS-AA, TraPPE-UA), formalize protocol adaptation into a rigorous decision framework, increase ensemble sizes~\cite{suter2025ensemble}, and extend validation to semicrystalline polymers, copolymers, and blends.



%%%%%%%%%%%%%%%%%%%%%%%%%%%%%%%%%%%%%%%%%%%%%%%%%%%%%%%%%%%%%%%%%%%%%
%% Technology Use Disclosure
%%%%%%%%%%%%%%%%%%%%%%%%%%%%%%%%%%%%%%%%%%%%%%%%%%%%%%%%%%%%%%%%%%%%%
\section*{Technology Use Disclosure}
The PolyJarvis system uses Anthropic's Claude large language model as the autonomous agent backbone.
Claude was also used to assist with grammar and typographical corrections during manuscript preparation.
All simulation results, analysis, and scientific conclusions have been carefully verified by the authors.


%%%%%%%%%%%%%%%%%%%%%%%%%%%%%%%%%%%%%%%%%%%%%%%%%%%%%%%%%%%%%%%%%%%%%
%% Code Availability
%%%%%%%%%%%%%%%%%%%%%%%%%%%%%%%%%%%%%%%%%%%%%%%%%%%%%%%%%%%%%%%%%%%%%
\section*{Data and Software Availability}
\texttt{https://github.com/arz-2/PolyJarvis}\\
Simulation trajectories and output files are available upon request.
\scaffold{R1 Major 12 / R1 Minor 12 --- Strengthen data \& software availability}{%
``Available upon request'' is inadequate for an automation paper. Replace (\placeholder{Zenodo DOI + release tag pending}):
\begin{itemize}
\item Deposit an archived release (cite exact version tag) --- not just a live GitHub link.
\item Include: system prompts / agent instructions, MCP tool schemas, all input LAMMPS scripts, random seeds, final structure files, analysis notebooks, and raw $T_g$-fit data for every benchmark system.
\item Provenance (workflow traces, tool calls, decision logs) is part of the core result, not supplementary --- make it openly available.
\item Update the statement to reference the archived-release DOI.
\end{itemize}}



%%%%%%%%%%%%%%%%%%%%%%%%%%%%%%%%%%%%%%%%%%%%%%%%%%%%%%%%%%%%%%%%%%%%%
%% Acknowledgments
%%%%%%%%%%%%%%%%%%%%%%%%%%%%%%%%%%%%%%%%%%%%%%%%%%%%%%%%%%%%%%%%%%%%%
\begin{acknowledgement}
Computational resources were provided by Lambda Labs.
The authors thank the RadonPy development team (Y.\ Hayashi, R.\ Yoshida, and collaborators) for making RadonPy available as open-source software.
\end{acknowledgement}


%%%%%%%%%%%%%%%%%%%%%%%%%%%%%%%%%%%%%%%%%%%%%%%%%%%%%%%%%%%%%%%%%%%%%
%% Author Declarations
%%%%%%%%%%%%%%%%%%%%%%%%%%%%%%%%%%%%%%%%%%%%%%%%%%%%%%%%%%%%%%%%%%%%%
\section*{Author Declarations}

\subsection*{Conflict of Interest}
The authors have no conflicts to disclose.

\subsection*{Author Contributions}
\textbf{Alexander Zhao:} Conceptualization, methodology, software development, simulation execution, data analysis, writing---original draft.
\textbf{Achuth Chandrasekhar:} Methodology, software development, writing---review and editing.
\textbf{Amir Barati Farimani:} Conceptualization, supervision, funding acquisition, writing---review and editing.


%%%%%%%%%%%%%%%%%%%%%%%%%%%%%%%%%%%%%%%%%%%%%%%%%%%%%%%%%%%%%%%%%%%%%
%% SI
%%%%%%%%%%%%%%%%%%%%%%%%%%%%%%%%%%%%%%%%%%%%%%%%%%%%%%%%%%%%%%%%%%%%%
\begin{suppinfo}
Section~S1: MCP tool inventory and server architecture;
Section~S2: detailed agent decision-making protocol;
Section S3: complete property calculation methods;
Section~S4: $T_g$ extraction methodology;
Section~S5: run-by-run $T_g$ summary;
Section~S6: system size and chain length;
Section~S7: bulk modulus volume statistics;
Section~S8: structural characterization;
Section~S9: complete per-run simulation parameters;
Section~S10: agent decision logs and autonomous error resolution summary;
Section~S11: equilibration convergence metrics;
Section~S12: computational resources and GPU cost summary.
\end{suppinfo}


%%%%%%%%%%%%%%%%%%%%%%%%%%%%%%%%%%%%%%%%%%%%%%%%%%%%%%%%%%%%%%%%%%%%%
%% Bibliography
%%%%%%%%%%%%%%%%%%%%%%%%%%%%%%%%%%%%%%%%%%%%%%%%%%%%%%%%%%%%%%%%%%%%%
\bibliography{reference.bib}

\end{document}
